% Generated. Do not edit.
% Projection scriptural-date-location/evidence, version 1: the Scriptural Date and
% Location page, generated whole from research/chronology.toml, the Scripture
% chronology corpus and the source library by tools/tpt document-types refresh.
% document: liturgy/roman-rite/1962/propers/temporal/54-fourteenth-after-pentecost
\section*{Appendix: Evidence for the Scriptural Dates}
\label{sec:date-location-evidence}
{\small For each passage, as the Scriptural Date and Location table names it: what is dated, how the corpus weighs each answer, the source's own words, and the sources cited. A subject the corpus does not date is reported as such and given no date.\par}
{\small
\par\smallskip\noindent\textit{Offertory: Ps. 33:8--9 (Vg.; Heb. 34)}\par\nobreak
\begin{itemize}[leftmargin=1.2em,itemsep=1pt,topsep=1pt,parsep=0pt]
\item \textbf{Traditional attribution} (disputed). David's reign in the usual chronology: B.C.~1055--1015 (Corbett also reports writers dating the reign 30--50 years later) (``reigned from 1055 to 1015 B.C.''; The Catholic Encyclopedia, Volume 4 (New York, 1908), article 'David, King,' printed p. 642, opening paragraph). Place: Hebron, then Jerusalem (The Holy Bible, Douay-Rheims (Challoner revision), Project Gutenberg eBook 1581: 2 Kings 5:5).
\item \textbf{Superscription setting} (preferred). David's flight to Achis, king of Geth: A.M.~2944 (Haydock edition reporting Ussher) (``A. M. 2944, A. C. 1060''; Douay Rheims Edition of the Holy Catholic Bible with a Comprehensive Catholic Commentary, William Peter Feeney 2014 Memorial Edition, 1K.21.1-15; Douay Rheims Edition of the Holy Catholic Bible with a Comprehensive Catholic Commentary, William Peter Feeney 2014 Memorial Edition, Ps.70.1). Place: Geth, at the Philistine court (The Holy Bible, Douay-Rheims (Challoner revision), Project Gutenberg eBook 1581: Ps. 33:1, 1 Kings 21:10, 21:15).
\item \textbf{Composition} (preferred). Psalms: Before c.~165 B.C. (``before the Maccabean period, around 165 B.C.; no individual psalm can be dated securely''; Catholic critical chronology; New American Bible, Revised Edition --- USCCB web state (2026), Psalms introduction).
\end{itemize}
\par\smallskip\noindent\textit{Introit: Ps. 83:10--11 (Vg.; Heb. 84), Ps. 83:2--3 (Vg.; Heb. 84)}\par\nobreak
\begin{itemize}[leftmargin=1.2em,itemsep=1pt,topsep=1pt,parsep=0pt]
\item \textbf{Historical setting}. No date: the chronology corpus has not yet dated this (research pending).
\item \textbf{Composition} (preferred). The Psalms of the Sons of Korah: Between the days of Isaias and the return from exile (The Catholic Encyclopedia, Volume 12 (1911)).
\end{itemize}
\par\smallskip\noindent\textit{Alleluia: Ps. 94:1 (Vg.; Heb. 95)}\par\nobreak
\begin{itemize}[leftmargin=1.2em,itemsep=1pt,topsep=1pt,parsep=0pt]
\item \textbf{Traditional attribution}. As under Offertory: Ps. 33:8--9 (Vg.; Heb. 34).
\item \textbf{Historical setting}. As under Introit: Ps. 83:10--11 (Vg.; Heb. 84), Ps. 83:2--3 (Vg.; Heb. 84).
\item \textbf{Composition}. As under Offertory: Ps. 33:8--9 (Vg.; Heb. 34).
\end{itemize}
\par\smallskip\noindent\textit{Gradual: Ps. 117:8--9 (Vg.; Heb. 118)}\par\nobreak
\begin{itemize}[leftmargin=1.2em,itemsep=1pt,topsep=1pt,parsep=0pt]
\item \textbf{Historical setting}. As under Introit: Ps. 83:10--11 (Vg.; Heb. 84), Ps. 83:2--3 (Vg.; Heb. 84).
\item \textbf{Composition}. As under Offertory: Ps. 33:8--9 (Vg.; Heb. 34).
\end{itemize}
\par\smallskip\noindent\textit{Gospel: Matt. 6:24--33}\par\nobreak
\begin{itemize}[leftmargin=1.2em,itemsep=1pt,topsep=1pt,parsep=0pt]
\item \textbf{Narrated event}. No date: the chronology corpus has not yet dated this (research pending).
\item \textbf{Composition} (disputed). The Gospel of St Matthew: c.~A.D.~38--45 (``about A.D. 38-45''; The Catholic Encyclopedia, Volume 10 (1911); The Catholic Encyclopedia, Volume 14 (New York, 1912)); c.~A.D.~40--42 (``about the year 40-42''; The Catholic Encyclopedia, Volume 10 (1911)); A.D.~40--45 (``the years 40-45''; The Catholic Encyclopedia, Volume 10 (1911)); c.~A.D.~60--68 (``about the year 60-68''; The Catholic Encyclopedia, Volume 10 (1911)); c.~A.D.~64--67 (``about the years 64-67''; The Catholic Encyclopedia, Volume 10 (1911)); c.~A.D.~50 (``about the year 50''; The Catholic Encyclopedia, Volume 14 (New York, 1912)). Place: Semitic original associated with Palestine (The Catholic Encyclopedia, Volume 10 (1911)).
\end{itemize}
\par\smallskip\noindent\textit{Communion: Matt. 6:33}\par\nobreak
\begin{itemize}[leftmargin=1.2em,itemsep=1pt,topsep=1pt,parsep=0pt]
\item \textbf{Composition}. As under Gospel: Matt. 6:24--33.
\end{itemize}
\par\smallskip\noindent\textit{Epistle: Gal. 5:16--24}\par\nobreak
\begin{itemize}[leftmargin=1.2em,itemsep=1pt,topsep=1pt,parsep=0pt]
\item \textbf{Composition} (answers differ). The Epistle to the Galatians: preferred: A.D.~58 (``anno Christi 58''; Commentaria in omnes D. Pauli epistolas (1614), prooemium; Commentaria in omnes D. Pauli epistolas (1614), epistola-ad-galatas-argumentum); alternative: A.D.~49--50 (``A.D. 49-50''; The Catholic Encyclopedia, Volume 6 (1909)); alternative: c.~A.D.~53--54 (``about A.D. 53, 54''; The Catholic Encyclopedia, Volume 6 (1909)); alternative: A.D.~56 (``(1 and 2 Corinthians; Galatians), 56''; The Catholic Encyclopedia, Volume 11 (New York, 1911)); alternative: A.D.~57--58 (``A.D. 57 or 58''; The Catholic Encyclopedia, Volume 6 (1909)).
\end{itemize}
\par}
